\section{Auxiliary structural reductions}\label{app:structural-auxiliary}
These arguments do not depend on the width-two proof. They concern only
the stated local laws and reductions, not extensions to larger width.

\begin{lemma}[Canonical pair laws on a forest]\label{lem:canonical-pairs}
On a forest of binary variables of prescribed means $p_i$, let $Q_{ij}$
be the midpoint of the two Fr\'echet endpoint pair laws on each edge:
\[
 Q_{ij}(1,1)=\frac{\min(p_i,p_j)+\max(0,p_i+p_j-1)}2.
\]
For any target set $S$ there is a global law with these pair marginals
and target union probability at least
$[\min(1,\sum_{i\in S}p_i)+\max_{i\in S}p_i]/2$, with both expressions
zero for $S=\varnothing$. For any nonnegative target payoff, such a law
also attains at least half its unconstrained singleton-feasible maximum.
The resulting law may depend on the target.
\end{lemma}
\begin{proof}
Every pair cell is affine in its $(1,1)$ probability and nonnegative
on the Fr\'echet interval, so its midpoint value is at least half every
feasible cell value. Choose a singleton-feasible law $P$ maximizing the
target union, using consecutive circular event intervals to attain
$c=\min(1,\sum_Sp_i)$. By the midpoint property, the edge measures
$R_{ij}=2Q_{ij}-P_{ij}$ are nonnegative probability measures with the
prescribed singleton means. Rooting each forest component and sampling
each child conditionally on its parent glues them into a law $R$ with
these edge marginals; null conditioning states are irrelevant. The
mixture $(P+R)/2$ has exactly the desired $Q$ marginals, and its target
union is at least $(c+b)/2$, since every law has union at least
$b=\max_Sp_i$. For a nonnegative payoff $g$, use the same construction
with a maximizing law $P$ for $g$ and $\E_Rg\ge0$.
\end{proof}
The residual construction fails on a triangle. For three fair bits,
let $P$ put mass $1/2$ on each of the two constant states. Each
canonical pair law is uniform, so $2Q_{ij}-P_{ij}$ requires perfect
anticorrelation, but three bits cannot be pairwise anticorrelated. This
is a failure of the residual-gluing step, not a counterexample to a
positive-monomial width-two bound.

\begin{lemma}[Identical-neighborhood compression]\label{lem:twin-compression}
In a positive unit-cube monomial factorization, let $B$ be a nonempty
variable block such that every factor containing a variable of $B$
contains all of $B$.
Replace it by one bit $W$ of mean
$w=\max(0,\sum_{i\in B}x_i-|B|+1)$ in those factors.
The full lower envelope and sum of local lower envelopes equal those
of the reduced instance. For some $\Delta\ge0$,
\[
 \tgap_{\rm original}=\tgap_{\rm reduced}+\Delta,\qquad
 \hgap_{\rm original}=\hgap_{\rm reduced}+\Delta.
\]
Thus every reduced gap inequality with constant at least one transfers.
\end{lemma}
\begin{proof}
For any original binary law, the product $V=\prod_{i\in B}X_i$ has
mean $q\ge w$. If $q>0$, independently retain its ones with probability
$w/q$, producing $W\le V$ of mean $w$; if $q=0$ use $W=0$.
The reduced objective is nondecreasing in $W$, so this operation cannot
increase it. Conversely, choose a local law on $B$ with all the original
means and product mean $w$, and attach it to a reduced law by
conditioning on its product value $W$. This lifts the reduced law to an
original law with the same objective and the original means. These two operations prove
equality of the full lower envelopes, including the endpoint cases.
For a factor with scope $B\cup C$, its original local lower value is
\[
 \max(0,\sum_Bx_i+\sum_Cx_i-|B|-|C|+1)
 =\max(0,w+\sum_Cx_i-|C|),
\]
where both sides vanish if $\sum_Bx_i-|B|+1<0$. The upper value can
only decrease under compression, because $w\le\min_Bx_i$; this also
holds when $C$ is empty and the reduced factor is affine.
Common upper attainment lets these decreases sum to one
$\Delta\ge0$. The two gap identities follow, and
$\tgap_{\rm reduced}+\Delta\le c(\hgap_{\rm reduced}+\Delta)$ for
$c\ge1$ proves the transfer without dividing by a gap.
\end{proof}

\begin{proposition}[Signed forest gluing on arbitrary boxes]
\label{prop:signed-forest}
For any multiaffine local factors on finite boxes whose incidence graph
is a forest, both upper and lower envelopes add. In particular this
holds for frequency-two factors whose dual multigraph is a forest.
\end{proposition}
\begin{proof}
Substitute fixed coordinates first. Choose an optimizing vertex law for
each factor and normalize its coordinate endpoints. Root each
incidence-tree component at a factor and traverse it, sampling each new
factor conditionally on its single already sampled variable. Consistent
singleton means make the laws agree at that separator; every other
variable of the factor is new, or there would be an incidence cycle.
Null separator states have zero probability. This preserves all chosen
optimizing laws simultaneously. Different components and unused variables
are sampled independently. Minimizing laws give the lower identity and
maximizing laws the upper identity. Signs do not enter the proof.
\end{proof}

For positive monomials, a partition of the factors into $q$ incidence
forests would give the bound $q$ by mixing the exact forest laws, but
treewidth does not bound the least such $q$. In $K_{2,m}$, $m\ge2$, with the two-vertex
side interpreted as variables, any two factors create a four-cycle;
all $m$ factors need distinct forest colors, although treewidth is two.
Adding one private variable to each factor makes their monomial scopes
distinct and preserves the obstruction. This is why the stronger
matrix argument is useful.

The general cutoff following Theorem~\ref{thm:variable-radix} also
shows a limitation of the fixed-radix construction. For integers
$b\ge2,L\ge2$ its profile capacities satisfy
\[
 M_s=\frac{(L-s)(b-1)+b}{b^s},\qquad
 \sum_{q=s+1}^L M_q=\frac{L-s}{b^s}.
\]
The sum identity follows by telescoping
$M_q=(L-q+1)b^{1-q}-(L-q)b^{-q}$.
If $s$ brackets one between $M_s$ and $M_{s+1}$, the exact hull is
$s+(L-s)b^{-s}$, as proved by its two attaining profiles.
For fixed $b$, $s=\log_b L+O_b(1)$: the inequality $M_s\ge1$ gives
$s\le\log_b((b-1)L+b)$, hence $L-s\sim L$; the reverse bracket then
gives $b^{s+1}\ge(b-1)(L-s-1)+b$. The remainder term is $O_b(1)$,
so the hull is $\log_b L+O_b(1)$. This limits only the fixed-radix
construction; it is not a sharp second-order lower theorem for all
positive polynomials.
